\section{Uncertainty and Additional Diagnostics}
\label{app:diagnostics}

The main text treats $H_{\mathrm{test}}$, $R_{\mathrm{abs}}$, and
$R_{\mathrm{KL}}$ as primary. Table~\ref{tab:metric-agreement} reports the
complementary $R_{\mathrm{NLL}}$ diagnostic,
Table~\ref{tab:question-level-similarity} gives paired-bootstrap intervals for
both twin-similarity ratios and $P_{\mathrm{closer}}$, and
Table~\ref{tab:nll-selectivity} reports the post-hoc NLL selectivity contrast.

\clearpage

\begin{table}[h]
  \centering
  \small
  \begin{tabular}{llrrr}
    \toprule
    Concept & Method & $H_{\mathrm{test}}\uparrow$ & $R_{\mathrm{NLL}}\downarrow$ & $R_{\mathrm{KL}}\downarrow$ \\
    \midrule
    Rome & EMBER & 0.208 & 1.365 & 2.264 \\
         & RMU   & \textbf{0.486} & \textbf{0.298} & 1.196 \\
         & SNMF  & 0.000 & 0.893 & \textbf{0.891} \\
         & RMU+E & 0.000 & 1.088 & 1.982 \\
         & SNMF+E & 0.000 & 1.327 & 2.181 \\
    \midrule
    Baseball & EMBER & \textbf{0.639} & 2.435 & 2.082 \\
             & RMU   & 0.239 & \textbf{0.655} & 1.569 \\
             & SNMF  & 0.129 & 0.910 & \textbf{0.962} \\
             & RMU+E & 0.465 & 2.131 & 2.015 \\
             & SNMF+E & 0.563 & 2.448 & 2.108 \\
    \midrule
    AI & EMBER & 0.753 & 2.731 & 3.310 \\
       & RMU   & 0.483 & \textbf{0.022} & 1.628 \\
       & SNMF  & 0.129 & 0.864 & \textbf{0.992} \\
       & RMU+E & \textbf{0.829} & 2.807 & 3.420 \\
       & SNMF+E & 0.816 & 2.750 & 3.354 \\
    \bottomrule
  \end{tabular}
  \caption{Signed-mean NLL diagnostic alongside the primary
  $H_{\mathrm{test}}$ and $R_{\mathrm{KL}}$ metrics. $R_{\mathrm{NLL}}<1$
  indicates greater signed-mean similarity to the twin than the full model,
  but opposing question-level differences can cancel. E abbreviates EMBER;
  bold marks the best erasure condition in each concept and column.}
  \label{tab:metric-agreement}
\end{table}

\begin{table}[h]
  \centering
  \small
  \begin{tabular}{llccc}
    \toprule
    Concept & Method & $R_{\mathrm{abs}}$ [95\% CI] $\downarrow$
      & $R_{\mathrm{KL}}$ [95\% CI] $\downarrow$
      & $P_{\mathrm{closer}}$ [95\% CI] $\uparrow$ \\
    \midrule
    Ancient Rome & EMBER & 1.613 [1.151, 2.244] & 2.264 [1.747, 2.970] & 0.40 [0.26, 0.54] \\
                 & RMU   & \textbf{0.660 [0.513, 0.827]} & 1.196 [1.037, 1.412] & \textbf{0.64 [0.50, 0.76]} \\
                 & SNMF  & 0.940 [0.891, 0.986] & \textbf{0.891 [0.838, 0.940]} & \textbf{0.64 [0.50, 0.76]} \\
                 & RMU+E & 1.384 [1.015, 1.873] & 1.982 [1.560, 2.549] & 0.44 [0.30, 0.58] \\
                 & SNMF+E & 1.558 [1.128, 2.143] & 2.181 [1.686, 2.851] & 0.42 [0.28, 0.56] \\
    \midrule
    Baseball     & EMBER & 2.128 [1.542, 2.978] & 2.082 [1.578, 2.770] & 0.26 [0.14, 0.38] \\
                 & RMU   & 1.227 [0.877, 1.680] & 1.569 [1.343, 1.823] & 0.46 [0.32, 0.60] \\
                 & SNMF  & \textbf{0.967 [0.888, 1.060]} & \textbf{0.962 [0.923, 1.003]} & \textbf{0.66 [0.54, 0.78]} \\
                 & RMU+E & 2.079 [1.510, 2.903] & 2.015 [1.537, 2.654] & 0.34 [0.22, 0.48] \\
                 & SNMF+E & 2.241 [1.658, 3.101] & 2.108 [1.610, 2.784] & 0.28 [0.16, 0.40] \\
    \midrule
    AI           & EMBER & 2.112 [1.402, 3.167] & 3.310 [2.384, 4.558] & 0.38 [0.24, 0.52] \\
                 & RMU   & 1.006 [0.754, 1.344] & 1.628 [1.434, 1.867] & 0.54 [0.40, 0.68] \\
                 & SNMF  & \textbf{0.927 [0.876, 0.976]} & \textbf{0.992 [0.954, 1.030]} & \textbf{0.64 [0.50, 0.76]} \\
                 & RMU+E & 2.196 [1.484, 3.251] & 3.420 [2.492, 4.644] & 0.34 [0.22, 0.48] \\
                 & SNMF+E & 2.152 [1.420, 3.239] & 3.354 [2.413, 4.609] & 0.40 [0.26, 0.54] \\
    \bottomrule
  \end{tabular}
  \caption{Target proximity to the concept-excluded twin.
  $R_{\mathrm{abs}}$ compares mean absolute correct-answer NLL differences,
  whereas $R_{\mathrm{KL}}$ compares mean full-vocabulary KL; values below one
  indicate improvement over the full model. $P_{\mathrm{closer}}$ is the
  fraction of questions on which the method has smaller absolute NLL distance.
  Intervals are paired-bootstrap percentile 95\% CIs over 50 questions (10,000
  resamples, seed 0). Bold marks the best erasure condition in each concept
  and column.}
  \label{tab:question-level-similarity}
\end{table}

Section~\ref{sec:results} discusses the main disagreements among these
diagnostics.

\clearpage

We used a post-hoc NLL contrast to summarize whether a method changed target
answers more than the two preservation groups. For method $M$ and concept
$c$, we defined
\begin{equation}
\begin{split}
S_{\mathrm{NLL}}(M,c)
&=\Delta\mathrm{NLL}^{\mathrm{target}}_{M-F}
-\frac{1}{2}\Delta\mathrm{NLL}^{\mathrm{neighbor}}_{M-F} \\
&\quad-\frac{1}{2}\Delta\mathrm{NLL}^{\mathrm{SciQ}}_{M-F}.
\end{split}
\end{equation}
A positive value indicates that support for correct target answers decreased
more than support for correct answers in the preservation groups. We also
report the two component contrasts so that the average does not conceal a
large change on either preservation group. To use one uncertainty procedure
for all five erasure conditions, we applied a normal approximation to every
row. Because the three test sets are disjoint, we estimated the standard error
from their across-question sample standard deviations,
\begin{equation}
\operatorname{SE}(S_{\mathrm{NLL}})
=\sqrt{\frac{s_T^2+\tfrac{1}{4}s_N^2+\tfrac{1}{4}s_S^2}{50}},
\end{equation}
and formed normal-approximation 95\% confidence intervals as
$S_{\mathrm{NLL}}\pm1.96\operatorname{SE}$. These intervals describe
variation across the fixed test questions for one trained checkpoint; they do
not capture variation across training seeds.

\begin{table}[h]
  \centering
  \footnotesize
  \renewcommand{\arraystretch}{1.05}
  \begin{tabular}{llrrr}
    \toprule
    Concept & Method & Target $-$ Neighbor & Target $-$ SciQ
      & $S_{\mathrm{NLL}}$ [95\% CI] \\
    \midrule
    Ancient Rome & EMBER & \textbf{1.468} & \textbf{2.395} & \textbf{1.932 $[1.343,\ 2.520]$} \\
                 & RMU   & $-0.306$ & 0.117 & $-0.094$ $[-0.442,\ 0.253]$ \\
                 & SNMF  & 0.087 & 0.174 & 0.131 $[0.070,\ 0.191]$ \\
                 & RMU+EMBER & 1.280 & 2.121 & 1.701 $[1.178,\ 2.224]$ \\
                 & SNMF+EMBER & 1.445 & 2.326 & 1.886 $[1.308,\ 2.463]$ \\
    \midrule
    Baseball     & EMBER & 1.194 & 1.289 & 1.242 $[0.700,\ 1.783]$ \\
                 & RMU   & 0.108 & 0.195 & 0.151 $[-0.239,\ 0.542]$ \\
                 & SNMF  & 0.127 & 0.049 & 0.088 $[-0.008,\ 0.184]$ \\
                 & RMU+EMBER & 1.092 & 1.159 & 1.126 $[0.569,\ 1.682]$ \\
                 & SNMF+EMBER & \textbf{1.229} & \textbf{1.343} & \textbf{1.286 $[0.732,\ 1.840]$} \\
    \midrule
    AI           & EMBER & 2.063 & 2.184 & 2.124 $[1.375,\ 2.872]$ \\
                 & RMU   & $-0.177$ & 0.199 & 0.011 $[-0.351,\ 0.373]$ \\
                 & SNMF  & 0.112 & 0.130 & 0.121 $[0.037,\ 0.205]$ \\
                 & RMU+EMBER & \textbf{2.128} & \textbf{2.523} & \textbf{2.325 $[1.580,\ 3.071]$} \\
                 & SNMF+EMBER & 2.109 & 2.182 & 2.145 $[1.374,\ 2.917]$ \\
    \bottomrule
  \end{tabular}
  \caption{Post-hoc target-specificity analysis on held-out questions. All
  entries are contrasts between signed NLL changes from the full model, in
  nats per answer token. Higher values indicate that the loss increase was
  more concentrated on target questions. The statistic does not measure
  proximity to the concept-excluded twin and was not used for checkpoint
  selection. All intervals use the normal approximation defined above.
  Bold marks the largest value in each concept and column.}
  \label{tab:nll-selectivity}
\end{table}
